%% Chen-Yu Kuo — resume (English)
%% Build: tectonic -X compile cv-en.tex
%%
%% Layout uses resume.cls by Weitian LI (LPPL 1.3c), vendored here verbatim:
%%   https://github.com/liweitianux/resume
%% itself based on YACC by Christophe Roger and Plasmati Graduate CV.

\documentclass{resume}

%% resume.cls only loads xeCJK under the [zh] option; this English document
%% still needs CJK glyphs for a handful of proper nouns.
%% Density tuning. resume.cls defaults to margin=1.5cm and \linespread{1.15};
%% this CV carries more content, so tighten a little without changing the look.
\usepackage{geometry}
\geometry{top=1.0cm, bottom=0.9cm, left=1.15cm, right=1.15cm, includefoot}
\linespread{1.0}
\setlength{\parskip}{0.35ex plus 0.3ex minus 0.1ex}
\setlist[itemize,1]{label=\faAngleRight, leftmargin=1.2em, labelsep=0.45em,
                    topsep=0.2ex, itemsep=0.25ex, parsep=0pt}
\titlespacing{\section}{0em}{0.6ex}{0.6ex}

%% Scale the class's fonts down a touch so the content lands on two pages.
\setmainfont{IBM Plex Serif}[Scale=0.96]
\setmonofont{IBM Plex Mono}[Scale=0.90]

\usepackage{xeCJK}
\setCJKmainfont{Noto Serif CJK TC}[Scale=0.96]
\XeTeXlinebreaklocale "zh"

\fileinfo{%
  \faGithub{} \link{https://github.com/whes1015/cv}{whes1015/cv} \hspace{0.5em}
  Layout: \link{https://github.com/liweitianux/resume}{liweitianux/resume} (LPPL~1.3c) \hspace{0.5em}
  \faEdit{} \today
}

\name{Chen-Yu}{Kuo}
\keywords{Go, TypeScript, Python, C++, Flutter, eBPF, embedded, seismology, earthquake early warning}
\tagline{郭宸毓 \enspace\textbullet\enspace Systems \& full-stack engineer --- earthquake early warning,
  from sensor firmware to the app that rings}

\profile{
  \email{whes1015@gmail.com}
  \github{whes1015}
  \iconlink{\faRobot}{https://huggingface.co/YuYu1015}{\texttt{huggingface.co/YuYu1015}} \\
  \degree{B.S. in Electrical Engineering}
  \university{National Kaohsiung Univ. of Science and Technology}
  \address{Kaohsiung, Taiwan}
}

\begin{document}
\makeheader

%======================================================================
\begin{abstract}
Founder and lead engineer of \textbf{ExpTech} (探索科技), a 29-person volunteer organisation running
Taiwan's largest independent earthquake early-warning platform. I work the whole path a seismic alert
travels --- ESP32 firmware written register by register, the multi-region alerting fleet it reports to, the
ground-motion models that decide what the alert says, and the Flutter app 100,000+ people installed.
\end{abstract}

%======================================================================
\sectionTitle{Competences}{\faWrench}
%======================================================================
\begin{competences}[9em]
  \comptence{Languages}{%
    TypeScript/JavaScript, Go, Python, Dart, C/C++, Rust, Lua, Shell
  }
  \comptence{Systems}{%
    Linux, eBPF/XDP, Docker, Nginx/OpenResty, CoreDNS, WebSocket/SSE, Prometheus, CI/CD
  }
  \comptence{Embedded}{%
    ESP32/Arduino, FreeRTOS, I\textsuperscript{2}C/SPI, LSM6DS3, GNSS (QZSS L1S), LoRa/Meshtastic
  }
  \comptence{Machine Learning}{%
    PyTorch, SeisBench, ONNX, scikit-learn, vLLM, llm-compressor, TensorRT Model Optimizer
  }
  \comptence{Domain}{%
    Earthquake early warning, GMPE calibration, CWA/JMA intensity scales, phase picking
  }
  \comptence{\icon{\faLanguage} Languages}{%
    \textbf{Mandarin} --- native; \textbf{English} --- technical reading \& writing
  }
\end{competences}

%======================================================================
\sectionTitle{Education}{\faGraduationCap}
%======================================================================
\begin{educations}
  %% TODO: replace [present] with your expected graduation, e.g. [Jun 2027]
  \education%
    {2022}%
    [present]%
    {National Kaohsiung Univ. of Science \& Technology}%
    {Dept. of Electrical Engineering}%
    {Electrical Engineering}%
    {B.S.}
\end{educations}

%======================================================================
\sectionTitle{Experience}{\faBriefcase}
%======================================================================
\begin{experiences}
  \experience%
    [Oct 2020]%
    {present}%
    {\textbf{Founder \& Lead Engineer} --- ExpTech (探索科技), Taiwan}%
    [\begin{itemize}
      \item Lead a 29-member volunteer organisation across 234 repositories as one of its two
        administrators, across firmware, 37 Go services, Flutter clients and ML pipelines. Authored
        \textbf{13,411} commits inside it (19,284 across GitHub, 477 PRs) and am top contributor on every
        flagship codebase --- DPIP 838/2,339, TREM-Lite 448/847, platform API 480/803 --- and sole author
        of the 665-commit production infrastructure repository.
      \item Signed a data agreement with Taiwan's \textbf{Central Weather Administration (CWA)} in late
        2022 --- reported by \emph{親子天下} as the Seismological Center's first non-corporate partner;
        invited back 10+ times as a collaborator (\emph{天下雜誌}, Feb 2026).
      \item Ship under real stakes: \textbf{100,000+} Google Play installs and 4.2/5 from 1,378
        App Store ratings for DPIP; \textbf{328,542} installer downloads across TREM desktop
        releases. During the M7.2 Hualien earthquake (3 Apr 2024) the network recorded
        800+ events in a single day.
    \end{itemize}]
\end{experiences}

%======================================================================
\sectionTitle{Machine Learning, Security \& Open Source}{\faFlask}
%======================================================================
\begin{itemize}
  \item \textbf{LLM inference \& model compression} --- 27 models on Hugging Face
    ({\raise0.06ex\hbox{$\sim$}}10,600 downloads). Quantise open models to \textbf{NVFP4} (TensorRT Model
    Optimizer), INT4-AutoRound and GGUF for vLLM and llama.cpp, up to 35B MoE. Implemented
    \textbf{refusal-direction ablation} --- deriving the refusal direction from paired prompts and
    orthogonalising it out of the weight matrices --- and released the \emph{Ornith-1.0} 9B/35B line
    including a DPO-tuned variant.
  \item \textbf{Security \& incident response} --- operated a honeypot that captured in-the-wild
    exploitation of \textbf{CVE-2025-55182} (React Server Components, CVSS 10.0) within days of its
    December 2025 disclosure; documented six malware samples with verified hashes and fifteen
    attacker source addresses.
  \item \textbf{Open source} --- 12 merged PRs to outside projects: a rendering-performance fix to the
    seismograph suite \textbf{AnyShake}
    (\link{https://github.com/anyshake/spectrogram-js}{\texttt{spectrogram-js}}), a
    Leaflet~$\rightarrow$~MapLibre migration for \texttt{smc.peering.tw}, and Flutter upgrades to
    \textbf{NKUST-AP-Flutter}, my university's student app.
  \item \textbf{Press} --- \emph{天下雜誌}, Feb 2026 (quoted as an ExpTech collaborator in a CWA
    Seismological Center profile); \emph{親子天下}, Jul 2024 (DPIP feature,
    {\raise0.06ex\hbox{$\sim$}}500,000 downloads after the 3 Apr 2024 quake).
\end{itemize}

\newpage
%======================================================================
\sectionTitle{Selected Engineering Work}{\faCode}
%======================================================================
\begin{projects}
  \project%
    {Mar 2026}[Jul 2026]%
    {sole author}%
    {Python, PyTorch, SeisBench, ONNX}%
    {Ground-motion \& phase-picking models}%
    {\begin{itemize}
      \item Refit Taiwan's shaking-intensity model on \textbf{45,398} CWA PGA/PGV/intensity observations
        (2000--2025), replacing a legacy equation that over-amplified with magnitude. Added a direct PGV
        prediction equation and corrected a site-amplification double-count that inflated estimates
        {\raise0.06ex\hbox{$\sim$}}2$\times$. Held-out: \textbf{62.0\% exact intensity vs 39.7\%} for the
        model then in production, MAE \textbf{0.400 vs 0.671}, 0\% false alarms.
      \item Trained a 50\,Hz PhaseNet P/S phase picker (6,457-line pipeline: dataset ingest, noise
        injection, evaluation, ONNX export). Diagnosed a windowing bug where the source window was
        shorter than the model window, teaching the model that P always sits mid-window and emitting a
        false pick every 45\,s stride. Both models published on Hugging Face.
    \end{itemize}}

  \project%
    {Sep 2025}[Jun 2026]%
    {sole author, 7,759 LOC}%
    {C++, ESP32, FreeRTOS}%
    {ES-Net --- seismometer firmware \& national sensor network}%
    {\begin{itemize}
      \item Hand-wrote an \textbf{LSM6DS3 accelerometer driver at register level} --- own register
        map, \textsc{who\_am\_i}, \textsc{ctrl1\_xl}, \textsc{int1\_ctrl} and FIFO paths --- rather
        than depending on a vendor library.
      \item On-device \textbf{JMA seismic-intensity} computation: an IIR filter bank plus a binary
        indexed tree for the running amplitude percentile the JMA scale requires.
      \item Timing is the hard problem here: SNTP in \textsc{smooth} mode so the clock slews rather
        than steps mid-waveform, plus HW/SW watchdogs, OTA update, and a documented custom binary wire
        protocol with TOTP station authentication.
      \item Feeds \textbf{200+ TREM-Net stations} nationwide; first-generation units were
        self-assembled for under NT\$800 each.
    \end{itemize}}

  \project%
    {Jan 2023}[present]%
    {top contributor, 448/847 commits}%
    {JavaScript, Electron, Rust}%
    {TREM --- real-time earthquake monitoring desktop platform}%
    {\begin{itemize}
      \item Desktop client rendering live ground motion from TREM-Net and CWA feeds; 93 GitHub stars and
        \textbf{328,542 installer downloads} across 74 releases, shipped continuously since Jan 2023.
        Designed the plugin system third parties extend it through and maintain its registry (417 commits);
        wrote the \texttt{trem-monitor} station-health service in Rust (105/150).
    \end{itemize}}

  \project%
    {Dec 2025}[present]%
    {sole author}%
    {Go, eBPF/XDP, OpenResty, CoreDNS}%
    {Production infrastructure}%
    {\begin{itemize}
      \item \textbf{kekkai} --- packet filter in Go using \texttt{cilium/ebpf}, with an XDP ingress
        program and a TC egress hook; \textsc{lru\_hash}/\textsc{array}/\textsc{ringbuf} maps drive
        per-IP rate limiting, allow/blocklists and port policy, with an emergency detach path.
      \item \textbf{Nginx/OpenResty fleet} --- 23 upstreams across nodes in Taipei, Kaohsiung,
        Tainan, Tokyo and Osaka, with Lua modules for traffic accounting, NTP and image proxying.
      \item \textbf{Split-horizon CoreDNS} with a custom Go \texttt{dnsstats} plugin
        (\texttt{miekg/dns}) exposing per-domain/client query stats and Prometheus metrics. Sole author
        too of the NTP server, API gateway and Meshtastic LoRa ingest service.
    \end{itemize}}

  \project%
    {Aug 2023}[present]%
    {top contributor, 838/2,339 commits}%
    {Dart/Flutter, TypeScript}%
    {DPIP --- cross-platform disaster alerting app}%
    {\begin{itemize}
      \item iOS and Android client integrating TREM-Net and CWA feeds with push alerting, MapLibre layers
        and 10 interface languages. Wrote the test suite and custom commit-gate tooling outright, and
        authored the backend service \texttt{dpip-server-ts} in full.
    \end{itemize}}

  \project%
    {Aug 2026}[present]%
    {sole author, {\raise0.06ex\hbox{$\sim$}}12k LOC with unit tests}%
    {C++}%
    {QZSS L1S disaster-message decoder}%
    {\begin{itemize}
      \item Decodes QZSS L1S 250-bit SBAS frames end to end --- preamble, 6-bit message type,
        212-bit payload and CRC-24Q --- including type 43 \textbf{DC Report} (JMA disaster
        broadcasts, IS-QZSS-DCR) and type 44 DCX.
    \end{itemize}}
\end{projects}


\end{document}
